\ifdefined\gkver\else\def\gkver{student}\fi
\documentclass[\gkver]{gaokaozhenti}
\begin{document}

\chapter{2006年天津理综}
%% paperType: 理综物理部分

\begin{enumerate}
\item
%% number: 14
%% paperName: 2006年天津理综第 14 题 6 分
%% typeId: 1
%% type: 单选题
%% chapter: 气体性质
%% point: 热力学第一定律
%% method: 无
%% score: 6
%% degree: 850
%% duplicateId: 0
%% body:
下列说法中正确的是\xzanswer{C}

\fourchoices[answer=C]
{任何物体的内能就是组成物体的所有分子热运动动能的总和}
{只要对内燃机不断改进，就可以把内燃机得到的全部内能转化为机械能}
{做功和热传递在改变内能的方式上是不同的}
{满足能量守恒定律的物理过程都能自发进行}

%% answer: C
%% memo:
\memoanswer{物体的内能是组成该物体的所有分子热运动的动能和分子势能的总和，故 A 错误；由热力学第二定律可知，不可能将内能全部用来对外做功而不引起其他变化，因此内燃机不可能将得到的内能全部转化为机械能，故 B 错误；做功是将其他形式的能转化为内能，而热传递则是物体间内能的转移，故 C 正确；满足能量守恒定律的物理过程有的具有方向性，故 D 错误。}

\item
%% number: 15
%% paperName: 2006年天津理综第 15 题 6 分
%% typeId: 1
%% type: 单选题
%% chapter: 光学
%% point: 光的折射
%% method: 无
%% score: 6
%% degree: 650
%% duplicateId: 0
%% body:
空气中两条光线 a 和 b 从方框左侧入射，分别从方框下方和上方射出，其框外光线如右图所示。方框内有两个折射率 $n=1.5$ 的玻璃全反射棱镜。下图给出了两棱镜四种放置方式的示意图，其中能产生右图效果的是\xzanswer{B}

\onepicture[w=3.2cm]{figs/15.png}

\fourchoices[answer=B, ispicture=true, h=2.2cm]
{figs/15a.png}
{figs/15b.png}
{figs/15c.png}
{figs/15d.png}

%% answer: B
%% memo:
\memoanswer{由棱镜的全反射规律作光路图可知，B 项正确。}

\item
%% number: 16
%% paperName: 2006年天津理综第 16 题 6 分
%% typeId: 1
%% type: 单选题
%% chapter: 曲线运动
%% point: 平抛运动
%% method: 无
%% score: 6
%% degree: 650
%% duplicateId: 0
%% body:
在平坦的垒球运动场上，击球手挥动球棒将垒球水平击出，垒球飞行一段时间后落地。若不计空气阻力，则\xzanswer{D}

\fourchoices[answer=D]
{垒球落地时瞬间速度的大小仅由初速度决定}
{垒球落地时瞬时速度的方向仅由击球点离地面的高度决定}
{垒球在空中运动的水平位移仅由初速度决定}
{垒球在空中运动的时间仅由击球点离地面的高度决定}

%% answer: D
%% memo:
\memoanswer{垒球的运动是平抛运动，由平抛运动规律可得，垒球落地速度为 $v_{1}=\sqrt{v_{0}^{2}+v_{y}^{2}}=\sqrt{v_{0}^{2}+2gh}$，其中 $h$ 为垒球下落的高度，故 A 错误；同理，垒球落地时速度的方向与水平方向的夹角 $\varphi=\arctan\frac{v_{y}}{v_{0}}=\arctan\frac{\sqrt{2gh}}{v_{0}}$，与高度 $h$ 有关，故 B 错误；水平位移 $x=v_{0}t=v_{0}\sqrt{\frac{2h}{g}}$，故 C 错误；而运动时间 $t=\sqrt{\frac{2h}{g}}$，故 D 正确。}

\item
%% number: 17
%% paperName: 2006年天津理综第 17 题 6 分
%% typeId: 1
%% type: 单选题
%% chapter: 机械振动
%% point: 振动图象
%% method: 无
%% score: 6
%% degree: 650
%% duplicateId: 0
%% body:
一单摆做小角度摆动，其振动图象如图，以下说法正确的是\xzanswer{D}

\onepicture[w=5.0cm]{\includesvg{figs/17}}

\fourchoices[answer=D]
{$t_{1}$ 时刻摆球速度最大，悬线对它的拉力最小}
{$t_{2}$ 时刻摆球速度为零，悬线对它的拉力最小}
{$t_{3}$ 时刻摆球速度为零，悬线对它的拉力最大}
{$t_{4}$ 时刻摆球速度最大，悬线对它的拉力最大}

%% answer: D
%% memo:
\memoanswer{由单摆振动规律可知，摆球在平衡位置时速度最大，做圆周运动的向心力最大，悬线对摆球的拉力最大。图象中 $t_{2}$、$t_{4}$ 时刻摆球在平衡位置，当摆球在最大位移处时速度最小，向心力为零，悬线拉力最小，$t_{1}$、$t_{3}$ 时刻摆球在最大位移处，故 D 正确。}

\item
%% number: 18
%% paperName: 2006年天津理综第 18 题 6 分
%% typeId: 1
%% type: 单选题
%% chapter: 原子和原子核
%% point: 核裂变
%% method: 无
%% score: 6
%% degree: 650
%% duplicateId: 0
%% body:
一个 $\ce{^{235}_{92}U}$ 原子核在中子的轰击下发生一种可能的裂变反应，其裂变方程为 $\ce{^{235}_{92}U + ^{1}_{0}n -> X + ^{94}_{38}Sr + 2^{1}_{0}n}$，则下列叙述正确的是\xzanswer{A}

\fourchoices[answer=A]
{X 原子核中含有 86 个中子}
{X 原子核中含有 141 个核子}
{因为裂变时释放能量，根据 $E=\Delta mc^{2}$，所以裂变后的总质量数增加}
{因为裂变时释放能量，出现质量亏损，所以生成物的总质量数减少}

%% answer: A
%% memo:
\memoanswer{在 $\ce{^{235}_{92}U}$ 的一种可能裂变方程中，由原子核反应中电荷数和质量数守恒可知，X 的核电荷数 $a=92-38=54$，质量数即核子数 $b=235+1-94-2=140$，故核内有 $140-54=86$ 个中子，54 个质子，故 A 正确、B 错误；由于裂变释放能量，有质量亏损，总质量会减小，但质量数不变，故 C、D 错误。}

\item
%% number: 19
%% paperName: 2006年天津理综第 19 题 6 分
%% typeId: 1
%% type: 单选题
%% chapter: 电路
%% point: 动态分析
%% method: 无
%% score: 6
%% degree: 450
%% duplicateId: 0
%% body:
如图所示的电路中，电池的电动势为 $E$，内阻为 $r$，电路中的电阻 $R_{1}$、$R_{2}$ 和 $R_{3}$ 的阻值都相同。在电键 S 处于闭合状态下，若将电键 $S_{1}$ 由位置 1 切换到位置 2，则\xzanswer{B}

\onepicture[w=4.5cm]{\includesvg{figs/19}}

\fourchoices[answer=B]
{电压表的示数变大}
{电池内部消耗的功率变大}
{电阻 $R_{2}$ 两端的电压变大}
{电池的效率变大}

%% answer: B
%% memo:
\memoanswer{开关 $S_{1}$ 由位置 1 切换到位置 2 后，电路总电阻 $R$ 减小，总电流 $I=\frac{E}{R}$ 增大，则电源内部消耗的功率 $P_{r}=I^{2}r$ 增大，故 B 正确；路端电压 $U=E-Ir$ 减小，因此电压表示数减小，故 A 错误；电阻 $R_{2}$ 两端电压由 $U_{2}=I_{2}R_{2}$ 变为 $U_{2}'=\frac{R_{3}}{R_{1}+R_{2}+R_{3}}I'R_{2}=\frac{1}{3}I'R_{2}$，因此 $U_{2}'<U_{2}$，故 C 错误；电池的效率为 $\eta=\frac{P_{\text{外}}}{P_{\text{外}}+P_{r}}\times100\%$，由于 $P_{\text{外}}$ 减小，$\eta$ 减小，故 D 错误。}

\item
%% number: 20
%% paperName: 2006年天津理综第 20 题 6 分
%% typeId: 1
%% type: 单选题
%% chapter: 电磁感应
%% point: 电磁感应中图象
%% method: 无
%% score: 6
%% degree: 450
%% duplicateId: 0
%% body:
在竖直向上的匀强磁场中，水平放置一个不变形的单匝金属圆线圈，规定线圈中感应电流的正方向如图 1 所示，当磁场的磁感应强度 $B$ 随时间 $t$ 如图 2 变化时，下图中正确表示线圈中感应电动势 $E$ 变化的是\xzanswer{A}

\twopicture[showlabel=false, h=3.0cm]
{figs/20a.png}
{figs/20b.png}

\fourchoices[answer=A, ispicture=true, h=2.0cm]
{figs/20c.png}
{figs/20d.png}
{figs/20e.png}
{figs/20f.png}

%% answer: A
%% memo:
\memoanswer{由法拉第电磁感应定律得，感应电动势 $E=n\frac{\Delta\varphi}{\Delta t}=n\frac{\Delta B}{\Delta t}S$，由图可知，$0\sim1\Us$ 内的感应电动势 $E_{1}$ 最大，且 $E_{1}=2E_{3}$，$1\sim3\Us$ 内 $E_{2}=0$。由楞次定律可以判断，$0\sim1\Us$ 内磁感应强度 $B$ 增大，感应电流的磁场与原磁场反向，故感应电流方向为图示正方向，故 A 正确。}

\item
%% number: 21
%% paperName: 2006年天津理综第 21 题 6 分
%% typeId: 1
%% type: 单选题
%% chapter: 电场
%% point: 带电粒子的直线运动
%% method: 建立模型
%% score: 6
%% degree: 250
%% duplicateId: 0
%% body:
在显像管的电子枪中，从炽热的金属丝不断放出的电子进入电压为 $U$ 的加速电场，设其初速度为零，经加速后形成横截面积为 $S$、电流为 $I$ 的电子束。已知电子的电量为 $e$、质量为 $m$，则在刚射出加速电场时，一小段长为 $\Delta l$ 的电子束内电子个数是\xzanswer{B}

\fourchoices[answer=B]
{$\frac{I\Delta l}{eS}\sqrt{\frac{m}{2eU}}$}
{$\frac{I\Delta l}{e}\sqrt{\frac{m}{2eU}}$}
{$\frac{I}{eS}\sqrt{\frac{m}{2eU}}$}
{$\frac{I\Delta lS}{e}\sqrt{\frac{m}{2eU}}$}

%% answer: B
%% memo:
\memoanswer{由题设模型有 $I=\frac{Q}{\Delta t}=\frac{ne}{\Delta t}=\frac{ne}{\Delta l}v$，电子在电场中加速时有 $eU=\frac{1}{2}mv^{2}$，联立得 $n=\frac{I\Delta l}{e}\sqrt{\frac{m}{2eU}}$，故 B 正确。}

\item
%% number: 22
%% paperName: 2006年天津理综第 22 题 4 分
%% typeId: 4
%% type: 实验题
%% chapter: 电路
%% point: 测电动势与内阻
%% method: 无
%% score: 4
%% degree: 650
%% duplicateId: 0
%% body:
\begin{enumerate}
	\item
	用半径相同的两小球 A、B 的碰撞验证动量守恒定律，实验装置示意如图，斜槽与水平槽圆滑连接。实验时先不放 B 球，使 A 球从斜槽上某一固定点 C 由静止滚下，落到位于水平地面的记录纸上留下痕迹。再把 B 球静置于水平槽前端边缘处，让 A 球仍从 C 处由静止滚下，A 球和 B 球碰撞后分别落在记录纸上留下各自的痕迹。记录纸上的 O 点是重锤线所指的位置，若测得各落点痕迹到 O 点的距离：OM $=2.68\Ucm$，OP $=8.62\Ucm$，ON $=11.50\Ucm$，并知 A、B 两球的质量比为 $2:1$，则未放 B 球时 A 球落地点是记录纸上的\tkanswer[1.5cm]{P}点，系统碰撞前总动量 $p$ 与碰撞后总动量 $p'$ 的百分误差 $\frac{|p-p'|}{p}=$\tkanswer[1.5cm]{2}\%（结果保留一位有效数字）。

	\onepicture[w=5.0cm, align=right]{figs/22a.png}

	\item
	一多用电表的电阻档有三个倍率，分别是 $\times 1$、$\times 10$、$\times 100$。用 $\times 10$ 档测量某电阻时，操作步骤正确，发现表头指针偏转角度很小，为了较准确地进行测量，应换到\tkanswer[1.8cm]{$\times 100$}档。如果换档后立即用表笔连接待测电阻进行读数，那么缺少的步骤是\tkanswer[2.8cm]{调零（或重新调零）}，若补上该步骤后测量，表盘的示数如图，则该电阻的阻值是\tkanswer[2.4cm]{$2.2\times 10^{3}\UO$}。

	\onepicture[w=8.0cm, align=right]{figs/22b.png}

	\item
	某研究性学习小组利用右图所示电路测量电池组的电动势 $E$ 和内阻 $r$。根据实验数据绘出如下图所示的 $R-\frac{1}{I}$ 图线，其中 $R$ 为电阻箱读数，$I$ 为电流表读数，由此可以得到 $E=$\tkanswer[1.5cm]{$2.9\UV$}，$r=$\tkanswer[1.5cm]{$0.9\UO$}。

	\twopicture[align=right, showlabel=false, hA=3.0cm, hB=4.5cm]
	{\includesvg{figs/22c}}
	{\includesvg{figs/22d}}
\end{enumerate}

%% answer:
\jdanswer{
\begin{enumerate}
	\item
	$P$；$2$

	\item
	$\times 100$；调零（或重新调零）；$2.2\times 10^{3}\UO$（或 $2.2\UkO$）

	\item
	$2.9\UV$；$0.9\UO$
\end{enumerate}
}

%% memo:
\memoanswer{
（1）A 与 B 相撞后，B 的速度增大，A 的速度减小，都做平抛运动，所以水平方向 B 在 A 的前面；小球离开水平槽后做平抛运动，A 与 B 下落的高度相同，在空中的运动时间相同，由于小球在水平方向上做匀速直线运动，小球的运动时间相等，因此小球的水平位移与小球的初速度成正比，计算时可以用小球的水平位移表示小球的初速度。根据题目所给实验数据，求出实验的百分误差。

A 与 B 相撞后，B 的速度增大，A 的速度减小，碰前碰后都做平抛运动，高度相同，落地时间相同，所以 P 点是没有碰时 A 球的落地点，N 是碰后 B 的落地点，M 是碰后 A 的落地点；系统碰撞前总动量 $p$ 与碰撞后总动量 $p'$ 的百分误差为

$\frac{|p-p'|}{p}=\frac{|m_{A}v_{A}-(m_{A}v_{A}'+m_{B}v_{B})|}{m_{A}v_{A}}=\frac{|2m_{B}\overline{OP}-(2m_{B}\overline{OM}+m_{B}\overline{ON})|}{2m_{B}\overline{OP}}\approx2\%$。

（2）用 $\times 10$ 挡测电阻时，表头偏转角度很小，说明待测电阻阻值很大，应换到 $\times 100$ 挡，换挡后应先欧姆调零，由表盘可知，其测量电阻的阻值为 $R=22.0\times100\UO=2.2\times10^{3}\UO$。

（3）由闭合电路的欧姆定律知 $I=\frac{E}{R+r}$，得 $\frac{1}{I}=\frac{R+r}{E}=\frac{1}{E}R+\frac{r}{E}$。由图象知，当 $\frac{1}{I}=1.0\UA^{-1}$ 时，$R=2\UO$；当 $\frac{1}{I}=3\UA^{-1}$ 时，$R=7.8\UO$，代入上式得 $E=2.9\UV$，$r=0.9\UO$。
}

\item
%% number: 23
%% paperName: 2006年天津理综第 23 题 16 分
%% typeId: 6
%% type: 计算题
%% chapter: 运动和力
%% point: 动量守恒定律
%% method: 无
%% score: 16
%% degree: 450
%% duplicateId: 0
%% body:
如图所示，坡道顶端距水平面高度为 $h$，质量为 $m_{1}$ 的小物块 A 从坡道顶端由静止滑下，进入水平面上的滑道时无机械能损失，为使 A 制动，将轻弹簧的一端固定在水平滑道延长线 M 处的墙上，一端与质量为 $m_{2}$ 的挡板 B 相连，弹簧处于原长时，B 恰位于滑道的末端 O 点。A 与 B 碰撞时间极短，碰后结合在一起共同压缩弹簧，已知在 OM 段 A、B 与水平面间的动摩擦因数均为 $\mu$，其余各处的摩擦不计，重力加速度为 $g$，求：

\begin{enumerate}
	\item
	物块 A 在与挡板 B 碰撞前瞬间速度 $v$ 的大小；
	\item
	弹簧最大压缩量为 $d$ 时的弹性势能 $E_{p}$（设弹簧处于原长时弹性势能为零）。
\end{enumerate}

\onepicture[w=6.0cm, align=right]{\includesvg{figs/23}}

%% answer:
\jdanswer{
\begin{enumerate}
	\item
	$v=\sqrt{2gh}$

	\item
	$E_{p}=\frac{m_{1}^{2}gh}{m_{1}+m_{2}}-\mu(m_{1}+m_{2})gd$
\end{enumerate}
}

%% memo:
\memoanswer{
（1）小物块运动的过程，由机械能守恒得 $m_{1}gh=\frac{1}{2}m_{1}v^{2}$，解得 $v=\sqrt{2gh}$。

（2）A、B 在碰撞过程中内力远大于外力，由动量守恒得 $m_{1}v=(m_{1}+m_{2})v'$；A、B 克服摩擦力所做的功为 $W=\mu(m_{1}+m_{2})gd$；由能量守恒定律得 $\frac{1}{2}(m_{1}+m_{2})v'^{2}=E_{\text{p}}+\mu(m_{1}+m_{2})gd$，解得 $E_{\text{p}}=\frac{m_{1}^{2}gh}{m_{1}+m_{2}}-\mu(m_{1}+m_{2})gd$。
}

\item
%% number: 24
%% paperName: 2006年天津理综第 24 题 18 分
%% typeId: 6
%% type: 计算题
%% chapter: 磁场
%% point: 洛伦兹力
%% method: 无
%% score: 18
%% degree: 450
%% duplicateId: 0
%% body:
在以坐标原点 $O$ 为圆心、半径为 $r$ 的圆形区域内，存在磁感应强度大小为 $B$、方向垂直于纸面向里的匀强磁场，如图所示。一个不计重力的带电粒子从磁场边界与 $x$ 轴的交点 A 处以速度 $v$ 沿 $-x$ 方向射入磁场，恰好从磁场边界与 $y$ 轴的交点 C 处沿 $+y$ 方向飞出。

\begin{enumerate}
	\item
	请判断该粒子带何种电荷，并求出其比荷 $\frac{q}{m}$；
	\item
	若磁场的方向和所在空间范围不变，而磁感应强度的大小变为 $B'$，该粒子仍从 A 处以相同的速度射入磁场，但飞出磁场时的速度方向相对于入射方向改变了 $60^{\circ}$ 角，求磁感应强度 $B'$ 多大？此次粒子在磁场中运动所用时间 $t$ 是多少？
\end{enumerate}

\onepicture[w=4.5cm, align=right]{\includesvg{figs/24}}

%% answer:
\jdanswer{
\begin{enumerate}
	\item
	$\frac{v}{Br}$

	\item
	$B'=\frac{\sqrt{3}}{3}B$，$t=\frac{\sqrt{3}\pi r}{3v}$
\end{enumerate}
}

%% memo:
\memoanswer{
（1）由粒子的飞行轨迹，利用左手定则可知，该粒子带负电荷。粒子由 A 点射入，由 C 点飞出，其速度方向改变了 $90^{\circ}$，即圆心角为 $90^{\circ}$，则粒子轨迹半径为 $R=r$，又 $qvB=m\frac{v^{2}}{R}$，则粒子的比荷 $\frac{q}{m}=\frac{v}{Br}$。

（2）粒子从 D 点飞出磁场速度方向改变了 $60^{\circ}$ 角，故 AD 弧所对圆心角为 $60^{\circ}$，粒子做圆周运动的半径为 $R'=r\cot30^{\circ}=\sqrt{3}r$，又 $R'=\frac{mv}{qB'}$，

\onepicture[w=4.5cm]{figs/24_an.jpg}

所以 $B'=\frac{\sqrt{3}}{3}B$。粒子在磁场中飞行时间为 $t=\frac{1}{6}T=\frac{1}{6}\times\frac{2\pi m}{qB'}=\frac{\sqrt{3}\pi r}{3v}$。
}

\item
%% number: 25
%% paperName: 2006年天津理综第 25 题 22 分
%% typeId: 6
%% type: 计算题
%% chapter: 曲线运动
%% point: 天体运动
%% method: 无
%% score: 22
%% degree: 250
%% duplicateId: 0
%% body:
神奇的黑洞是近代引力理论所预言的一种特殊天体，探寻黑洞的方案之一是观测双星系统的运动规律。天文学家观测河外星系大麦哲伦云时，发现了 LMCX$-$3 双星系统，它由可见星 A 和不可见的暗星 B 构成。两星视为质点，不考虑其它天体的影响，A、B 围绕两者连线上的 O 点做匀速圆周运动，它们之间的距离保持不变，如图所示。引力常量为 $G$，由观测能够得到可见星 A 的速率 $v$ 和运行周期 $T$。

\begin{enumerate}
	\item
	可见星 A 所受暗星 B 的引力 $F_{A}$ 可等效为位于 O 点处质量为 $m'$ 的星体（视为质点）对它的引力，设 A 和 B 的质量分别为 $m_{1}$、$m_{2}$，试求 $m'$（用 $m_{1}$、$m_{2}$ 表示）；
	\item
	求暗星 B 的质量 $m_{2}$ 与可见星 A 的速率 $v$、运行周期 $T$ 和质量 $m_{1}$ 之间的关系式；
	\item
	恒星演化到末期，如果其质量大于太阳质量 $m_{s}$ 的 2 倍，它将有可能成为黑洞。若可见星 A 的速率 $v=2.7\times10^{5}\Ums$，运行周期 $T=4.7\pi\times10^{4}\Us$，质量 $m_{1}=6m_{s}$，试通过估算来判断暗星 B 有可能是黑洞吗？（$G=6.67\times10^{-11}\UNmqkgq$，$m_{s}=2.0\times10^{30}\Ukg$）
\end{enumerate}

\onepicture[w=5.0cm, align=right]{\includesvg{figs/25}}

%% answer:
\jdanswer{
\begin{enumerate}
	\item
	$m'=\frac{m_{2}^{3}}{(m_{1}+m_{2})^{2}}$

	\item
	$m'=\frac{m_{2}^{3}}{(m_{1}+m_{2})^{2}}=\frac{v^{3}T}{2\pi G}$

	\item
	暗星 B 有可能是黑洞
\end{enumerate}
}

%% memo:
\memoanswer{
（1）设 A、B 的圆轨道半径分别为 $r_{1}$、$r_{2}$，A、B 之间的距离为 $r$，由题意知，A、B 做匀速圆周运动的角速度相同，设其为 $\omega$，由牛顿运动定律得 $F_{A}=m_{1}\omega^{2}r_{1}$，$F_{B}=m_{2}\omega^{2}r_{2}$，又 $F_{A}=F_{B}$，$r=r_{1}+r_{2}$，由上述各式得 $r=\frac{m_{1}+m_{2}}{m_{2}}r_{1}$。由万有引力定律得 $F_{A}=G\frac{m_{1}m_{2}}{r^{2}}$，将上式代入得 $F_{A}=G\frac{m_{1}m_{2}^{3}}{(m_{1}+m_{2})^{2}r_{1}^{2}}$，令 $F_{A}=G\frac{m_{1}m'}{r_{1}^{2}}$，比较可得 $m'=\frac{m_{2}^{3}}{(m_{1}+m_{2})^{2}}$。

（2）由牛顿第二定律得 $G\frac{m_{1}m'}{r_{1}^{2}}=m_{1}\frac{v^{2}}{r_{1}}$，又可见星 A 的轨道半径为 $r_{1}=\frac{vT}{2\pi}$，由以上各式解得 $\frac{m_{2}^{3}}{(m_{1}+m_{2})^{2}}=\frac{v^{3}T}{2\pi G}$。

（3）将 $m_{1}=6m_{s}$ 代入上式，得 $\frac{m_{2}^{3}}{(6m_{s}+m_{2})^{2}}=\frac{v^{3}T}{2\pi G}$，代入数据得 $\frac{m_{2}^{3}}{(6m_{s}+m_{2})^{2}}=3.5m_{s}$。设 $m_{2}=nm_{s}$（$n>0$），将其代入上式，得 $\frac{m_{2}^{3}}{(6m_{s}+m_{2})^{2}}=\frac{n}{(\frac{6}{n}+1)^{2}}m_{s}=3.5m_{s}$。可见 $\frac{m_{2}^{3}}{(6m_{s}+m_{2})^{2}}$ 的值随 $n$ 的增大而增大，试令 $n=2$，得 $\frac{n}{(\frac{6}{n}+1)^{2}}m_{s}=0.125m_{s}<3.5m_{s}$，若使上式成立，则 $n$ 必大于 2，即暗星 B 的质量 $m_{2}$ 必大于 $2m_{s}$，由此得出结论：暗星 B 有可能是黑洞。
}

\end{enumerate}

\end{document}
